\paragraph{范数降低不等于因果误差分摊。}
对中档RFM的速度最大误差而言，延后换步恢复了相对恒步增量的约 $5.584\%$。这只是算术比值 $(E_8-E_{12})/(E_8-E_U)$，并未隔离桥接单步的误差。该干预同时改变了四个时间单位内的细步/粗步积分、桥接时刻与状态，以及后续误差历史。且 $t=8$ 位于解析脉冲内部，不是非光滑激励拼接点。

保存流场可直接说明非可加性。令 $e_8=u_8-u_{\rm ref}$、$e_{12}=u_{12}-u_{\rm ref}$、$\delta=e_8-e_{12}$，统一使用物理面积内积，并除以 $\|u_{\rm ref}\|^2$，则
\begin{equation}
 \|e_8\|_{\rm ref}^2=\|e_{12}\|_{\rm ref}^2+
 2\langle e_{12},\delta\rangle_{\rm ref}+\|\delta\|_{\rm ref}^2.
 \label{eq:review-switch-attribution}
\end{equation}
对RFM37种子7331，取 $t=72$、layer144网格及时间细化参考，四项依次为 $1.9956194141\times10^{-7}$、$1.9202943907\times10^{-7}$、$7.4475738974\times10^{-9}$、$8.4928450335\times10^{-11}$，均为未乘百分比尺度的相对误差平方。未舍入恒等式残差为 $2.04\times10^{-22}$；交叉项约占平方范数差的 $98.87\%$。这条确定性恒等式说明，$5.6\%$ 的恢复比例及其 $94.4\%$ 补数都不是独立的物理误差成分。若要因果区分桥接与后续粗步误差，仍需单独控制相应干预并做时间步细化研究。
